% ============================================================
% T2 超对称、扭化与拓扑场论 —— B 组 14 条
% ============================================================
\track{T2\quad 超对称、扭化与拓扑场论}{本轨道 14 条：扭化是把物理理论变成拓扑场论的手术，也是「物理侧能翻译成数学」的入口。}

\begin{word}{$N=4$ 超对称杨--米尔斯}{$N=4$ supersymmetric Yang--Mills}{n4-sym}{§1 p.3；§2.1 pp.8--11}
\mathdef{取定紧单群 $G$ 与四维流形 $M$。四维 $\mathcal N=4$ 超对称杨--米尔斯理论由以下数据组成：$G$-主丛 $E$ 上的联络 $A$（曲率 $F$），以及六个取值于伴随丛 $\mathrm{ad}E$ 的标量场 $\phi_i$，$i=0,\dots,5$；作用量为 $\frac{1}{e^2}\int\mathrm{Tr}\,(\frac12F\wedge\ast F+\sum_iD\phi_i\wedge\ast D\phi_i+\frac12\sum_{i,j}[\phi_i,\phi_j]^2)+\frac{\theta}{8\pi^2}\int\mathrm{Tr}\,F\wedge F$。它是超共形场论：超对称代数含 16 个超荷，复耦合 $\tau=\theta/2\pi+4\pi i/e^2$ 是精确边缘参数。}
\mathex{给定 $G$ 时它是四维唯一带最大超对称的规范理论。在平坦 $\mathbb R^4$ 上真空模空间为 $(\mathbb R^6\otimes\mathfrak t)/W$，$\mathfrak t$ 为 Cartan 子代数、$W$ 为 Weyl 群；$G=SU(2)$ 时它是 $\mathbb R^6/\mathbb Z_2$，一般点处 $SU(2)$ 破缺为 $U(1)$。}
\physdef{在 KW06 中它是全书的唯一输入：四维最大超对称规范理论，含规范场 $A_\mu$、六个标量 $\phi_i$ 与四个 Weyl 旋量超荷 $\lambda^i$。它本身不是拓扑场论；选取扭化后，它才变成一族以 $\mathbb{CP}^1$ 为参数的拓扑场论（§3）。}
\physex{取 $G=SU(2)$、$\theta=0$、$e\ll1$，则 $\tau=4\pi i/e^2$ 为纯虚。在 Coulomb 分支取 $\langle\phi\rangle=\mathrm{diag}(a,-a)$、$a\ne0$，则 $SU(2)\to U(1)$，W 玻色子质量为 $|Z|=\sqrt{2/\mathrm{Im}\,\tau}\,|a|$（式 (2.22) 取 $n=1,m=0$）。}
\bridge{物理侧这一个理论按 §3.1 的三种方式扭化，给出三种拓扑场论；其中第三种（GL 扭化）是几何 Langlands 的物理载体。它的 $S$-对偶 $\tau\mapsto-1/\tau$ 与扭化参数 $t$ 耦合，最终进入典范参数 $\Psi=\frac{\tau+\bar\tau}{2}+\frac{\tau-\bar\tau}{2}\frac{t-t^{-1}}{t+t^{-1}}$ 的 $SL(2,\mathbb Z)$ 作用（§3.5，式 (3.50)）。}
\intuit{它是四维规范理论的「氢原子」：超对称与共形对称把作用量锁死到只剩 $e$ 与 $\theta$，因此许多量可精确计算、对偶可精确检验。类比的边界是它没有引力，不是弦论。}
\pitfall{两个常见错：其一，把它当成 $\mathcal N=4$ 超引力——后者含引力子，本文的理论不含；其二，把 $\mathcal N=4$ 读成「4 个超荷」——它其实是 4 个最小超对称的拷贝，共 16 个超荷。}
\end{word}

\begin{word}{超荷}{supercharge}{supercharge}{§2.1 p.9}
\mathdef{超荷是超对称代数中的奇生成元 $Q_a$：把玻色子变为费米子，在四维按 $Spin(4)$ 的旋量表示变换，并与动量满足 $\{Q_a,Q_b\}=\sum_I\Gamma^I_{ab}P_I$（式 (2.9)）。数学上它们构成超 Poincaré 代数的奇部——一个 $\mathbb Z_2$ 分次李超代数，偶部是 Poincaré 代数；超荷是守恒流 $J^I=\frac12\mathrm{Tr}\,\Gamma^{JK}F_{JK}\Gamma^I\lambda$ 的荷。}
\mathex{四维 $\mathcal N=4$ 共有 16 个实超荷：约化后为四个 Weyl 旋量 $Q^i_\alpha$（$i=1,2,3,4$，$\alpha=1,2$），按 $SU(4)_{\mathcal R}$ 的 $\mathbf 4$ 变换，复共轭按 $\bar{\mathbf 4}$。$\mathcal N=1$ 只有一个 Weyl 旋量、4 个实超荷；超对称的「个数」按实分量计，不按旋量个数计。}
\physdef{在 §2.1 中，超荷来自十维超杨--米尔斯的 16 分量 Majorana--Weyl 旋量。在扭化理论里，超对称变分写成 $\delta_T\Phi=[Q,\Phi\}$，其参数 $\epsilon$ 取值于正旋量空间 $\mathcal S^+$，就代表一个超荷；能否在弯曲 $M$ 上定义超对称，取决于能否选到满足不变性条件的 $\epsilon$。}
\physex{$G=SU(2)$、$M=\mathbb R^4$：四个超荷 $\lambda^i$ 的变分为 $\delta\lambda=\frac12\sum F_{IJ}\Gamma^{IJ}\epsilon$（§3.1）。取常数 $\epsilon$ 且在扭化后的转动下不变，就得到拓扑超荷 $Q$；$Q$ 的变分不含导数，故与度规无关。}
\bridge{扭化直接作用在超荷上：$\epsilon\in\mathcal S^+$ 在 $Spin(4)\times SU(4)_{\mathcal R}$ 的表示中分解，取 $Spin'(4)$ 不变的组合即得标量超荷。GL 扭化保留 $\epsilon_\ell,\epsilon_r$ 两个独立解，故 $Q=uQ_\ell+vQ_r$，$t=v/u$；这正是一切拓扑场论的起点（§3.1）。}
\intuit{超荷是「平移的平方根」：$\{Q,Q\}=P$ 使能量非负、使 BPS 态的质量被中心荷锁住。扭化后这个平方根可以不伴随平移而单独存在，于是 $Q^2=0$——这是拓扑场论的关键。类比边界：普通超荷在一般弯曲流形上没有全局定义，标量超荷才有。}
\pitfall{「$\mathcal N=4$ 就是 4 个超荷」是错的，实为 16 个；「超荷与共形超荷是一回事」也错，后者是 $S_\alpha,\bar S_{\dot\alpha}$，与 $Q$ 共同组成超共形代数。KW06 扭化用的 $Q$ 是超荷与 R-对称的混合，不是普通 Poincaré 超荷。}
\end{word}

\begin{word}{标量超荷}{scalar supercharge}{scalar-supercharge}{§3.1 p.18}
\mathdef{标量超荷是扭化理论中在转动群下为标量、且满足 $Q^2=0$ 的超荷。精确地说：取同态 $\varkappa:Spin(4)\to Spin(6)$，用对角子群 $Spin'(4)=(1\times\varkappa)(Spin(4))$ 重定义转动；若 $\mathcal S^+$ 含非零的 $Spin'(4)$-不变旋量 $\epsilon$，则 $Q=\epsilon^aQ_a$ 是转动不变的奇算子，且 $\epsilon^T\Gamma^\mu\epsilon=0$ 给出 $Q^2=\frac12\{Q,Q\}=0$（模规范变换）。}
\mathex{数学原型是微分：$Q$ 是分次代数 $(A^\bullet,Q)$ 上的奇微分，$Q^2=0$，$H^\bullet(Q)=\ker Q/\mathrm{im}\,Q$。de Rham 微分 $d$、Dolbeault 微分 $\bar\partial$、Chevalley--Eilenberg 微分都是这种「标量奇算子」；扭化把它从几何搬到规范理论，$Q$-上同调于是成了拓扑不变量。}
\physdef{§3.1：扭化要求 $\mathcal S^+$ 中有 $Spin'(4)$ 不变向量，对应的超对称生成元就是标量超荷；它使路径积分可以只对 $Q$-不变场计算。原文常称它为拓扑超荷，§6.4 又把它称为 BRST 算子——同一个对象，三种叫法。}
\physex{$G=SU(2)$、$M=\mathbb R^4$：三种扭化中前两种各留一个不变旋量，GL 扭化留两个。取 $t=i$ 时 $Q=uQ_\ell+vQ_r$；$Q$-不变玻色场满足 $\mathcal F=0$，其中 $\mathcal A=A+i\phi$ 为复化联络（§3.2），即 $Q$-闭场恰是平坦 $G_{\mathbb C}$ 联络。}
\bridge{标量超荷是「物理侧能翻译成数学」的入口：$Q$ 是微分，$Q$-上同调是上同调。物理侧的 $Q$-闭算子、$Q$-不变路径积分，对应数学侧的上同调类与示性类。它也是拓扑场论定义的核心——能量动量张量是 $Q$-恰当的。}
\intuit{「标量」指它在扭化后的转动下不变：自旋指标被 R-对称「吃掉」了。直觉上像给旋量丛选了一个整体平凡化，于是旋量变成函数、超对称变成普通微分。边界：它只在扭化后才存在，未扭化理论没有这个标量超荷。}
\pitfall{把「标量超荷」当成普通超荷的某个分量，或以为它天然存在：它只在扭化后才有意义。而且「标量」是对 $Spin'(4)$ 而言，不是对原来的 $Spin(4)$。另一个错误是以为 $Q^2=0$ 普遍成立——未扭化时 $\{Q,Q\}=P\ne0$。}
\end{word}

\begin{word}{BRST 算子}{BRST operator}{brst-op}{§6.4 p.97}
\mathdef{数学上，BRST 算子是微分分次代数 $(A^\bullet,d)$ 的微分：$d^2=0$、$d$ 为奇导子。标准来源是李代数作用——$\mathfrak g$ 作用于复形 $V$ 时，Chevalley--Eilenberg 微分 $d_{CE}$ 给出复形 $(C^\bullet(\mathfrak g,V),d_{CE})$，其上同调 $H^\bullet(\mathfrak g,V)$ 即规范不变量的代数模型。一般地，BRST 上同调就是「奇算子模去其像」后的上同调。}
\mathex{$\mathfrak g=\mathfrak{su}(2)$、$V=\mathbb C$ 为平凡模：$H^0(\mathfrak{su}(2),\mathbb C)=\mathbb C$，$H^1=H^2=0$，$H^3=\mathbb C$（紧半单李代数的上同调是外代数，只在 $0$ 与 $3$ 次非零）。这个例子说明 BRST 上同调确实在算「不变量」，并没有把一切平凡化。}
\physdef{§6.4 把扭化理论中的拓扑超荷 $Q$ 直接称为 BRST 算子，并把 $Q$-恰当项称为 BRST-平凡项（如 §6.3 的 $\{Q,\cdot\}$）。物理上 $Q$ 的作用是「模掉规范」，但它不是规范固定意义下的 BRST 算子，而是超对称的守恒荷；区别在于后者依赖规范选取，前者是扭化的产物。}
\physex{$G=SU(2)$、$M=\Sigma\times C$、$t=\pm i$：§6.3 用 $Q$-恰当性说明，移动线算子的插入点 $p\in C$ 只把作用量改变一个 BRST-平凡项，故关联函数局部地与 $p$ 无关。这是 BRST 语言在 KW06 中最实际的用法。}
\bridge{「BRST 算子 = 标量超荷 = 拓扑超荷」是同一对象在三种语境下的名字：数学侧叫微分，物理侧叫超荷，规范理论传统叫 BRST。词典的规则是：看到 BRST 就找它的微分分次结构，看到 $Q$ 就取上同调。}
\intuit{BRST 算子像「模掉规范变换」的代数版本：闭合理的是物理态，恰当的是纯规范。直觉之所以靠得住，是因为 $Q^2=0$ 保证 $\mathrm{im}\,Q\subset\ker Q$，上同调良定义。边界：KW06 的 $Q$ 来自超对称，而非 Faddeev--Popov 规范固定。}
\pitfall{把 KW06 的 BRST 算子等同于通常的规范固定 BRST 算子：后者作用在鬼场与物理态之间、依赖规范选取；前者是扭化后的超荷，规范不变。另一个错是以为 BRST 上同调平凡——上同调非零正是拓扑不变量的来源。}
\end{word}

\begin{word}{$Q$-上同调}{cohomology of $Q$}{q-cohom}{§3.1 p.18}
\mathdef{设 $Q$ 是分次向量空间 $A^\bullet$ 上的奇算子、$Q^2=0$。$Q$-上同调定义为 $H^\bullet(Q)=\ker Q/\mathrm{im}\,Q$：$Q$-闭模去 $Q$-恰当。若 $A^\bullet$ 还是分次代数且 $Q$ 是导子，则 $H^\bullet(Q)$ 继承分次代数结构。典型例子：de Rham 复形的 $H^\bullet_{dR}(M)$（$Q=d$）、Dolbeault 复形的 $H^{p,q}_{\bar\partial}(X)$、李代数上同调 $H^\bullet(\mathfrak g,V)$。}
\mathex{$M=S^2$、$Q=d$：$H^0=\mathbb R$、$H^1=0$、$H^2=\mathbb R$，上同调维数等于 Betti 数；这是「$Q$-上同调算拓扑」的最简模型。扭化理论中把 $Q$ 换成标量超荷，$H^\bullet(Q)$ 的结构就成为四维流形的拓扑不变量。}
\physdef{§3.1：选定标量超荷 $Q$ 后，物理诠释改变——只关心 $Q$-不变的路径积分、算子与态，把 $[Q,\mathcal O\}$ 视为平凡；「有趣的算子落在 $Q$ 的适当上同调中」。态空间取 $Q$-上同调，$Q$-恰当态视为零。}
\physex{$G=SU(2)$、$M=\mathbb R^4$、$t=\pm i$：$Q$-不变玻色场是平坦 $G_{\mathbb C}$ 联络 $\mathcal A=A+i\phi$（§3.2），故 $Q$-上同调的「基态」由平坦联络的模空间参数化；$t$ 为实数时，$Q$-不变场改由椭圆方程 (3.29) 的解给出。}
\bridge{$Q$-上同调是词典的翻译机制：物理侧的 $Q$-上同调类对应数学侧的上同调类。例如 Donaldson 理论中 $Q$-上同调给出 Donaldson 不变量；在 KW06 中，Wilson/'t Hooft 线算子的 $Q$-上同调类对应 $C$ 上的 Hecke 特征层（§11，由 T5 轨道展开）。}
\intuit{像 de Rham 上同调「闭形式模恰当形式」：不是在数闭形式，而是在数「洞」。$Q$-上同调数的是扭化后剩下的拓扑自由度；之所以靠得住，是因为 $Q^2=0$ 使商空间良定义、且与度规无关。}
\pitfall{最常见的错是只取 $Q$-不变子空间而漏掉「模去 $Q$-恰当」：那样得到的是闭算子而非上同调类，维数会错。另一个错是把 $Q$-上同调当成「BRST 平凡化」，以为取上同调后什么都不剩；正相反，非平凡类（瞬子数、线算子）正是物理内容。}
\end{word}

\begin{word}{R-对称}{R-symmetry}{r-sym}{§2.1 p.10}
\mathdef{R-对称是超对称代数的自同构群中与 Poincaré 代数交换、却非平凡地作用在奇生成元（超荷）上的部分。对四维 $\mathcal N=4$，它是 $SU(4)_{\mathcal R}\cong Spin(6)$：四个超荷按 $\mathbf 4$（共轭按 $\bar{\mathbf 4}$）变换，六个标量按 $\mathbf 6=\wedge^2\mathbf 4$ 变换，即 $SO(6)$ 的矢量表示。数学上它是超共形代数 $\mathfrak{psu}(2,2|4)$ 的偶部中与 $\mathfrak{sl}(2)\oplus\mathfrak{sl}(2)$ 交换的 $\mathfrak{su}(4)$ 部分。}
\mathex{$SU(4)_{\mathcal R}$ 的中心由 $\mathcal I$ 生成，它在 $\mathbf 4$ 上作用为 $i$、在 $\bar{\mathbf 4}$ 上作用为 $-i$（§2.2）。子群 $SU(2)_\ell\times SU(2)_r\subset SU(4)_{\mathcal R}$ 使 $\mathbf 4=(\mathbf 2,\mathbf 1)\oplus(\mathbf 1,\mathbf 2)$；另有与它交换的 $U(1)$，生成元 $\mathcal K=i\,\mathrm{diag}(1,-1)$，$\mathbf 4$ 按 $\mathcal K$ 本征值 $+1,-1$ 分块。这套表示论是扭化的全部原料。}
\physdef{§2.1：十维理论约化到四维时，内部 $SO(6)$ 转动成为 R-对称 $SU(4)_{\mathcal R}$；它是整体对称性，与规范群 $G$ 交换，作用在四个超荷与六个标量上。扭化正是把这个 R-对称的子群混入 Lorentz 群，从而改变「自旋」的定义。}
\physex{$G=SU(2)$、$M=\mathbb R^4$：六个标量 $\phi_0,\dots,\phi_5$ 按 $\mathbf 6_v$ 变换；取 $\sigma=(\phi_4-i\phi_5)/\sqrt2$，则 $\mathcal K(\sigma)=2$、$\mathcal K(\bar\sigma)=-2$，而规范场 $A$ 的 $\mathcal K=0$、费米子 $\psi,\tilde\psi$ 的 $\mathcal K=1$。按 $\mathcal K$ 分次是 §3 全部符号簿记的骨架。}
\bridge{扭化的定义就是取同态 $\varkappa:Spin(4)\to Spin(6)=SU(4)_{\mathcal R}$，用 $Spin'(4)=(1\times\varkappa)(Spin(4))$ 替换原转动群。R-对称的表示分解（$\mathbf 4$、$\mathbf 6_v$、$\mathcal K$ 本征值）决定了扭化后各场的自旋与分次，是「物理侧怎么变成数学侧」的字典第一页。}
\intuit{R-对称是「转动超荷」的内禀对称性：像味对称，但它作用在超对称本身而不只作用在物质场上。之所以靠得住，是因为超对称代数允许这类自同构；扭化就是让这个内禀转动与时空转动「绑在一起」。边界：它是整体对称，不是规范对称。}
\pitfall{把 $SU(4)_{\mathcal R}$ 与规范群 $SU(4)$ 混淆；或以为扭化改变了理论的 R-对称——扭化不改变理论，只改变我们用哪个对角子群来定义「自旋」。另一个错是以为 R-对称只作用在费米子上：六个标量也按 $\mathbf 6$ 变换。}
\end{word}

\begin{word}{扭化}{twisting}{twisting}{§3.1 p.17}
\mathdef{扭化是把超对称理论变成拓扑场论的手术。设理论在 $M$ 上的对称群为 $Spin(4)\times G_R$，取同态 $\varkappa:Spin(4)\to G_R$，令 $Spin'(4)=(1\times\varkappa)(Spin(4))$ 为对角子群。扭化后的理论仍由同样的场组成，但场的自旋按 $Spin'(4)$ 而非 $Spin(4)$ 定义——等价于把场丛沿 $\varkappa$ 拉回，把 R-对称指标混入自旋指标。要求 $\mathcal S^+$ 含非零的 $Spin'(4)$-不变旋量，于是得到标量超荷 $Q$ 且 $Q^2=0$。}
\mathex{取 $G=SU(2)$、$M=\Sigma\times C$、$C$ 亏格 $g\ge2$：扭化后的四维理论紧化到 $\Sigma$ 上，得到靶为 Hitchin 模空间 $\mathcal M_H(SU(2),C)$ 的二维 sigma 模型，$\dim_{\mathbb C}\mathcal M_H=6(g-1)$，其 Hitchin 纤维化的底为 $H^0(C,K^2)$（$3(g-1)$ 维）。}
\physdef{§3.1：三种不等价的 $\varkappa$（差 Weyl 群作用）给出三种扭化；第三种（GL 扭化）是本文主角。扭化不是简化或截断：作用量没变，变的是「转动」的解释与只保留 $Q$-不变量的约定；代价是只在欧氏号差下好用。}
\physex{$G=SU(2)$、$M=\mathbb R^4$：扭化后玻色场重新分配——规范场仍是 1-形式，原六个标量中四个合成第二个 1-形式 $\phi$，另两个合成复标量 $\sigma$；费米子变成两个 1-形式、两个 2-形式 $\chi^\pm$ 与两个 0-形式 $\eta,\tilde\eta$。$Q$ 的变分不含被变场的导数（如 $\delta_TA_\mu=iu\psi_\mu+iv\tilde\psi_\mu$），这正是拓扑性的来源。}
\bridge{扭化是两个独立事实的合成：其一，超荷按 R-对称的某个子群（这里是 $SU(2)_\ell\times SU(2)_r\subset SU(4)_{\mathcal R}$）重新定义自旋，从而让标量超荷存在；其二，扭化后存在标量超荷 $Q$ 使 $Q^2=0$，其上同调给出拓扑不变量。前者让 $Q$ 存在，后者让 $Q$ 有用，缺一不可。}
\intuit{扭化像「换一副眼镜看同一个理论」：把 R-转动与时空转动绑在一起，带旋量指标的荷便看似标量。靠得住是因为这只是重组表示、不改变路径积分，改变的只是我们只问 $Q$-不变的东西。边界：只在欧氏号差下成立。}
\pitfall{「扭化是简化/截断」是头号误解：理论没变，变的是对称群的嵌入与诠释。另一个误解是「扭化后理论就平凡了」：$Q$-上同调非平凡，拓扑不变量正住在这里。也别与「加大超对称」混为一谈。}
\end{word}

\begin{word}{GL 扭化}{GL twist}{gl-twist}{§3.1 p.17}
\mathdef{GL 扭化是三种扭化中的第三种：$\varkappa:Spin(4)\to Spin(6)$ 的像在 $\mathcal S^+$ 上保留两个线性独立的不变旋量 $\epsilon_\ell,\epsilon_r$，故不变超荷构成二维复空间；整体标度无意义，留下 $\mathbb{CP}^1$ 参数族 $Q=uQ_\ell+vQ_r$，$t=v/u$。它是 Vafa--Witten 三种扭化中唯一保留两个超荷的一种，对应 Hitchin 模空间超凯勒结构的 twistor 族。}
\mathex{$\mathcal M_H(G)$ 有复结构 $I,J,K$；参数 $t$ 走遍与 twistor 族相应的 $\mathbb{CP}^1$：$t=\pm i$ 时进入复结构 $J$ 的 B-模型，实 $t$ 时进入辛形式 $\omega=\mathrm{Im}\,\Omega$ 的 A-模型。$S$-对偶把 $t$ 按 $t\mapsto\frac{c\tau+d}{|c\tau+d|}t$ 变换（式 (3.13)），在二维图像里就是 $\mathcal M_H(G)$ 与 $\mathcal M_H({}^LG)$ 的镜像对称。}
\physdef{§3.1：判据是 $\mathcal S^+$ 中不变旋量的个数：前两种扭化各一个，GL 扭化两个，故有二维复空间。$SU(4)_{\mathcal R}$ 的中心元 $\mathcal I$ 使 $\epsilon_\ell\to i\epsilon_\ell$、$\epsilon_r\to i^{-1}\epsilon_r$，从而 $t\to-t$；这是唯一的平凡等价（式 (3.14)）。}
\physex{$G=SU(2)$、$M=\mathbb R^4$：$t=\pm i$ 时 $Q$-不变场满足 $\mathcal F=0$，即 $\mathcal A=A+i\phi$ 平坦，给出同态 $\pi_1(M)\to G_{\mathbb C}$ 的模空间 $\mathcal Y_M$；实 $t$（如 $t=1$）时方程变成椭圆型，按解的带符号计数；$t=0,\infty$ 端点分别只剩瞬子、反瞬子解。$t$ 与 $\Psi$ 的换算：$t=\pm1$ 时 $\Psi=\mathrm{Re}\,\tau$，$t=\pm i$ 时 $\Psi=\infty$（§3.5）。}
\bridge{GL 扭化是几何 Langlands 的物理载体：实 $t$ 侧是 A-模型（靶 $\mathcal M_H$ 的辛结构），$t=\pm i$ 侧是 B-模型（靶为平坦 $G_{\mathbb C}$ 丛的复结构）；$S$-对偶对应镜像对称，$\Sigma\times C$ 紧化把 4d 理论变成 2d sigma 模型。这一条是 T2 通往 T3/T6 的门。}
\intuit{前两种扭化各只留下一个超荷，拓扑性很强但结构较死；GL 扭化留下两个，于是保留了一整个 $\mathbb{CP}^1$ 的超荷家族，足以在二维约化后看到完整的超凯勒几何。边界：多一个参数不等于多一个物理自由度——$t$ 只是理论族的一个参数。}
\pitfall{把 GL 扭化与 Donaldson--Witten 扭化混为一谈：后者是 $\mathcal N=2$ 的扭化，给出 Donaldson 不变量；GL 扭化是 $\mathcal N=4$ 的第三种扭化，通往几何 Langlands 的镜像对称。另一个错是以为 $t$ 是物理耦合：$t$ 与 $\tau$ 只有组合 $\Psi$ 有物理意义。}
\end{word}

\begin{word}{拓扑场论}{topological field theory}{tqft}{§1 p.2；§3.1 p.18}
\mathdef{拓扑场论在数学上的标准定义是 Atiyah--Segal 函子：$Z:\mathrm{Bord}_d\to\mathrm{Vect}$，把闭 $(d-1)$ 维流形映为向量空间、$d$ 维配边映为线性映射，并保持粘合与不交并。等价地，所有关联函数与度规无关。上同调场论版本：存在奇算子 $Q$ 使 $Q^2=0$、能量动量张量 $T_{\mu\nu}$ 是 $Q$-恰当的，可观测量取 $Q$-上同调。}
\mathex{二维 TQFT 等价于交换 Frobenius 代数：亏格 $g$ 闭曲面的配分函数是乘积的迹。四维例子是 Donaldson 理论，其观测量由瞬子模空间上的积分定义、是光滑四维流形的 Donaldson 不变量；三维例子是 Chern--Simons 理论。}
\physdef{§1、§3.1：扭化后的 $\mathcal N=4$ 理论「松散地说是拓扑场论」——原文特意加「松散地」，因为它可能依赖 $M$ 的光滑结构而不只依赖同伦型。判据是：度规改变只让作用量差一个 $Q$-恰当项，故配分函数与关联函数不随度规改变。}
\physex{$G=SU(2)$、$M$ 闭四维流形、$t$ 实：配分函数等于超对称方程 (3.29) 的解的带符号计数，符号为单圈费米子行列式的符号；计数在椭圆方程族的连续形变下不变，故配分函数与 $t$ 无关（§3.4）。$t=1$ 时这是本文的「A-模型」端。}
\bridge{拓扑场论是扭化的输出、也是词典的靶：物理侧的 TQFT 对应数学侧的不变量理论（Donaldson 不变量、Gromov--Witten 不变量、Hitchin 系统的计数）。KW06 的路线是 4d TQFT 紧化到 2d TQFT，再取 $Q$-上同调得到特征层。}
\intuit{TQFT 像「只记得拓扑的物理」：一台把流形切开再粘回去、输出数与向量空间的机器。之所以靠得住，是因为 $T_{\mu\nu}=\{Q,\cdot\}$ 意味着度规的连续形变不改变任何关联函数。边界：它仍可能记得光滑结构，且需要 $Q$ 存在。}
\pitfall{把「拓扑场论」理解成「没有度规的场论」：中间仍有度规，只是依赖都是 $Q$-恰当的。也不要把原文的「松散地说是 TQFT」当成严格定理：光滑结构依赖是真实的可能性。也别以为所有 TQFT 都来自扭化。}
\end{word}

\begin{word}{瞬子数}{instanton number}{inst-number}{§1 p.3；§2.2 p.12；§3.3 p.22}
\mathdef{设 $E$ 是闭可定向四维流形 $M$ 上的 $G$-丛，$F$ 为联络曲率。瞬子数定义为 $k=\frac{1}{8\pi^2}\int_M\mathrm{Tr}\,F\wedge F$（本文的反厄米约定），它等于 $E$ 的第二陈数 $\langle c_2(E),[M]\rangle$，取整数值，因此只依赖丛的拓扑。它是 $H^4(M,\mathbb Z)$ 中示性类的积分；方向反转时变号。}
\mathex{$G=SU(2)$、$M=S^4$：$G$-丛由 $k\in\mathbb Z\cong\pi_3(SU(2))$ 分类，$k=1$ 的 BPST 瞬子是显式解，其模空间维数为 $8k-\dim G=5$。$M$ 换成 $S^2\times S^2$ 时 $k$ 仍是整数，瞬子按两个 $S^2$ 上的陈数分解。}
\physdef{§1、§2.2：瞬子数是 $\theta$ 项所乘的拓扑荷，$\mathcal I_\theta=-\frac{\theta}{8\pi^2}\int\mathrm{Tr}\,F\wedge F$；瞬子作用量为 $8\pi^2|k|/e^2$，第 $k$ 扇区的权重为 $e^{i\theta k-8\pi^2k/e^2}$。§3.3 消灭定理：若 $\int\mathrm{Tr}\,F\wedge F\ne0$，则对 $t\ne0,\infty$ 无超对称解；$t=0,\infty$ 时瞬子、反瞬子解仍存在。}
\physex{$G=SU(2)$、$M=S^4$、$k=1$：BPST 瞬子的作用量为 $8\pi^2/e^2$，弱耦合 $e\ll1$ 时贡献 $e^{-8\pi^2/e^2}$，是典型的非微扰效应。定向反转使 $k\to-k$，从而 $\Psi\to-\Psi$（$t\to1/t$，§3.5）。}
\bridge{瞬子数把物理侧的拓扑荷翻译成数学侧的特征数：$k=\langle c_2(E),[M]\rangle$。它是 $\theta$ 角周期性的根源（整性使 $\theta\cong\theta+2\pi$），也是扭化理论消灭定理的核心；在紧化 $\Sigma\times C$ 上，它控制哪些超对称位形存在。}
\intuit{瞬子数像「缠绕数」：规范场绕时空「卷」了几圈。之所以靠得住，是因为它是示性类的积分，在连续形变下不变。边界：它数的是整张流形上的拓扑荷，不是「瞬子个数」那样的粒子计数。}
\pitfall{把瞬子数与瞬子作用量混同：作用量是 $8\pi^2|k|/e^2$，含耦合常数。另一个错是忽略定向反转使 $k\to-k$，漏掉符号会得到错误的 $T$ 变换。还有：对一般 $G$，取迹用哪个表示、归一化如何取，决定整性是否成立。}
\end{word}

\begin{word}{BPS 条件}{BPS condition}{bps}{§5.4 p.73}
\mathdef{BPS 条件是超对称理论中「质量达到中心荷下界」的条件：若超对称代数带中心荷 $Z$，则 $\{Q,Q\}$ 的半正定性给出 $M\ge|Z|$；等号成立、且部分超荷湮灭该态时，称该态为 BPS 态。数学上这意味着超对称表示变短（短多重态），态的数目少于一般多重态；BPS 态的稳定性、墙穿越与计数是 Hitchin 系统与计数几何的核心对象。}
\mathex{$\mathcal N=2$、$G=SU(2)$ 的 Seiberg--Witten 理论：电荷 $(n,m)$ 的 BPS 态质量 $M=|Z|$，$Z=an+a_Dm$，其中 $a$ 是标量期望值、$a_D=\partial\mathcal F/\partial a$；$Z$ 也是 Seiberg--Witten 微分 $\lambda$ 沿 $H_1(\Sigma,\mathbb Z)$ 上闭链的周期积分。$a_D=0$ 处单极子变无质量，$a=0$ 处 W 玻色子变无质量。}
\physdef{§5.4：在 KW06 中 BPS 条件就是「保住部分超对称」的条件——对粒子是 $M=|Z|$，对线算子与场位形则要求所规定的奇异性与拓扑方程相容。扭化后，BPS 方程正是 $Q$-不变场方程，例如 §3.2 的方程 (3.29)。}
\physex{$G=SU(2)$、$\theta=0$、$\tau=4\pi i/e^2$：电荷 $(1,0)$ 的 W 玻色子满足 $M=|Z|=\sqrt{2/\mathrm{Im}\,\tau}\,|a|$；电荷 $(0,1)$ 的单极子满足 $M=|Z|=\sqrt{2/\mathrm{Im}\,\tau}\,|a\tau|$。半 BPS 只保一半超荷；$\mathcal N=4$ 理论里还存在 1/4 BPS 位形。}
\bridge{BPS 条件把物理的「质量等于中心荷」翻译成数学的「稳定/自对偶方程」：Hitchin 方程、单极子方程、瞬子方程都是 BPS 方程；BPS 态的计数对应模空间上的相交理论。扭化后 $Q$-不变方程就是 BPS 方程，这是 T2 与 T3 的接口。}
\intuit{像测地线：在所有连接两点的曲线中，测地线是长度的临界点；BPS 态在其电荷扇区里是能量下界的饱和者。之所以靠得住，是因为下界来自 $\{Q,Q\}$ 的正定性，不依赖相互作用细节。边界：并非所有稳定态都 BPS，也并非所有 BPS 态都稳定。}
\pitfall{「BPS = 质量为零」是错的：只有中心荷为零的态才可能无质量，一般 BPS 态质量 $=|Z|\ne0$。另一个错是「BPS = 保持全部超对称」：半 BPS 只保一半。还有：不要把 BPS 条件与「基态」或「稳定态」混同。}
\end{word}

\begin{word}{中心荷}{central charge}{central-charge}{§2.2 p.14}
\mathdef{中心荷是超对称代数的中心元 $Z$：与所有生成元交换，出现在超荷反对易子的右端，如 $\{Q^i_{\dot A},Q^j_{\dot B}\}=\epsilon_{\dot A\dot B}\epsilon^{ij}Z$（式 (2.18)）。数学上它是超李代数的中心扩张；$\{Q,Q\}$ 的半正定性给出 $M\ge|Z|$，故中心荷是电荷格上的线性泛函，BPS 态的质量恰为其模。}
\mathex{$\mathcal N=2$、$G=SU(2)$：$Z=an+a_Dm$，$(n,m)\in\mathbb Z^2$ 为电、磁荷；它是 Seiberg--Witten 曲线 $\Sigma$ 上的微分 $\lambda$ 沿闭链 $\gamma$ 的周期 $Z=\oint_\gamma\lambda$，故中心荷的取值格与曲线的周期格直接相关。}
\physdef{§2.2：$\mathcal N=4$ 理论约化时被丢弃的六个动量 $P_{i+4}$ 的期望值（标量场在无穷远的渐近值）作为中心荷出现，它们就是电、磁荷。$S$-对偶按 $Z\to\frac{|c\tau+d|}{c\tau+d}Z$ 变换（式 (2.24)），相位是模 $2\pi$ 的；超荷是它的平方根，故拿到半相位、拼成双重覆盖的表示。}
\physex{$G=SU(2)$、$M=\mathbb R^4$：$Z(\vec n,\vec m)=\sqrt{2/\mathrm{Im}\,\tau}\,\vec\phi\cdot(\vec n+\tau\vec m)$（式 (2.22)）；取 $\phi=a$、$(n,m)=(1,0)$ 得 $Z=\sqrt{2/\mathrm{Im}\,\tau}\,a$。$S$-对偶交换电与磁：$(n,m)\to(m,-n)$，同时 $\tau\to-1/\tau$。}
\bridge{中心荷是 BPS 质量的物理来源，也是数学侧周期积分的物理对应：物理的电荷格对应数学的 $H_1(\Sigma,\mathbb Z)$，物理的 $Z$ 对应数学的 $\oint_\gamma\lambda$。KW06 用中心荷在 Coulomb 分支上算出 $S$-对偶对超荷的作用（式 (2.24) 开方得半相位，涉及 $(-1)^F=\mathcal I^2$）。}
\intuit{中心荷是「电荷在真空上的投影」：给定态的电荷向量与真空标量期望值，就得到一个数 $Z$，BPS 态的质量恰是这个数。像功等于力乘位移，$Z$ 是电荷与真空势的配对。边界：$Z$ 依赖真空（$a$ 的取值），不是常数。}
\pitfall{与二维共形场论的中心荷 $c$ 混淆：后者是 Virasoro 代数的中心项，与超对称中心荷毫无关系。另一个错是把中心荷当成规范荷：它是整体对称的荷，且同一电荷扇区的 $Z$ 在真空模空间上连续变化。}
\end{word}

\begin{word}{拓扑超荷}{topological supercharge}{top-supercharge}{§6.4 p.102}
\mathdef{拓扑超荷是标量超荷 $Q$ 的别名，强调它生成「拓扑对称」：$Q$ 是分次代数上的奇微分，$Q^2=0$，且对任意场 $\mathcal O$，$[Q,\mathcal O]$ 不含被变场的导数——因此度规改变只让作用量差一个 $Q$-恰当项。数学上它把拓扑场论定义中的微分具体化：$H^\bullet(Q)$ 就是拓扑可观测量的代数。}
\mathex{de Rham 微分是原型：$Q=d$、$Q^2=0$，$H^\bullet(Q)=H^\bullet_{dR}(M)$，可观测量是调和形式。扭化理论中 $Q$ 换成超荷，$H^\bullet(Q)$ 换成四维流形的拓扑不变量（如 Donaldson 不变量、Vafa--Witten 不变量）。}
\physdef{§6.4：KW06 用「拓扑超荷」称 $Q$；§6.3 又把 $Q$-恰当项叫「BRST-平凡项」。同一对象三个名字：标量超荷（强调自旋）、拓扑超荷（强调对称性）、BRST 算子（强调上同调）。它的存在是扭化理论成为拓扑场论的关键条件。}
\physex{$G=SU(2)$、$M=\mathbb R^4$、$t=i$：拓扑超荷 $Q=uQ_\ell+vQ_r$ 的变分为 $\delta_TA_\mu=iu\psi_\mu+iv\tilde\psi_\mu$、$\delta_T\phi_\mu=iv\psi_\mu-iu\tilde\psi_\mu$（§3.1）——右端不含导数，故 $Q$ 与度规无关；$Q$-不变场是平坦 $G_{\mathbb C}$ 联络。}
\bridge{拓扑超荷是词典的「动词」：它把物理侧的对称性变成数学侧的微分，把物理的不变量变成数学的上同调类。对应规则是：找到 $Q$（由扭化给），取 $H^\bullet(Q)$（由代数给），再识别为已知的数学上同调（de Rham、Dolbeault、相交上同调）。}
\intuit{「拓扑」指它看不见度规：像 $d$ 只依赖光滑结构，不依赖度量。之所以靠得住，是因为 $Q$ 的变分不含导数，故 $T_{\mu\nu}=\{Q,\cdot\}$，度规形变是 $Q$-恰当的。边界：这种拓扑性只在扭化后成立。}
\pitfall{把拓扑超荷与普通超荷当成两个不同对象：在 KW06 中它就是扭化后的那个标量超荷；也不要把「拓扑对称」理解成理论没有费米子或没有度规——费米子仍在，只是成对抵消。另一个错是以为拓扑超荷唯一：GL 扭化有 $\mathbb{CP}^1$ 族。}
\end{word}

\begin{word}{Coulomb 分支}{Coulomb branch}{coulomb-branch}{§2.2 p.14}
\mathdef{Coulomb 分支是超对称规范理论真空模空间的一个分支：伴随标量场取非零期望值 $\langle\phi\rangle\in\mathfrak t$（$\mathfrak t$ 为 Cartan 子代数），把 $G$ 破缺到极大环面 $U(1)^r$，$r=\mathrm{rank}\,G$。模去 Weyl 群 $W$ 后，分支局部同构于 $\mathfrak t_{\mathbb C}/W$；其坐标环是 $W$-不变多项式，如 $u=\langle\mathrm{Tr}\,\phi^2\rangle$。}
\mathex{$G=SU(2)$、$\mathcal N=2$：Coulomb 分支是 $u$-平面，$u=\langle\mathrm{Tr}\,\phi^2\rangle\in\mathbb C$，奇异点 $u=\pm\Lambda^2$ 处单极子或双子变无质量。$\mathcal N=4$ 时六个标量都可取期望值，真空模空间为 $\mathbb R^6/\mathbb Z_2$，其上度规平坦、$\tau$ 不跑动。}
\physdef{§2.2：KW06 在 Coulomb 分支上计算 $S$-对偶对超荷的作用——那里规范群破缺为 $U(1)^r$，可显式算出电、磁荷与中心荷，从而确定 $\hat\phi(\gamma)$ 与 $(-1)^F$ 的关系；低能阿贝尔理论的 $SL(2,\mathbb Z)$ 对偶也在此登场。}
\physex{$G=SU(2)$、$M=\mathbb R^4$：取 $\langle\phi\rangle=\mathrm{diag}(a,-a)$、$a\ne0$，W 玻色子质量为 $|Z|=\sqrt{2/\mathrm{Im}\,\tau}\,|a|$；$a\to-a$ 是 Weyl 反射，给出同一个真空。$a=0$ 处 $SU(2)$ 未破缺，是分支的原点。}
\bridge{Coulomb 分支是「物理侧能算」的地方：中心荷、BPS 质量、$S$-对偶作用都在这里显式计算，再借超对称与全纯性延拓到整个理论（§2.2--2.3）。数学对应是 Seiberg--Witten 曲线与周期积分；但注意 KW06 的 $\mathcal M_H$ 不是 Coulomb 分支。}
\intuit{Coulomb 分支像「指南针选了一个方向」：标量场在 Lie 代数里挑一个方向，连续族给出连续真空。之所以靠得住，是因为 $F$ 项平坦要求 $[\phi_i,\phi_j]=0$，可以同时对角化。边界：$\mathcal N=4$ 时它没有量子修正，$\mathcal N=2$ 时才有。}
\pitfall{头号错误是把 Coulomb 分支与 Hitchin 模空间 $\mathcal M_H$ 混为一谈：后者是二维 sigma 模型的靶、Higgs 丛的模空间，不是四维真空模空间。另一个错是与 Higgs 分支混同：后者是味对称破缺的分支，$\mathcal N=2$ 时二者不同。}
\end{word}

\par\medskip
\noindent{\small 本轨道小结：超对称给出超荷与 R-对称；扭化把 R-对称混入自旋，得到标量超荷 $Q$；$Q$-上同调把理论变成拓扑场论并提供翻译机制；瞬子数、中心荷、BPS 条件提供可算的拓扑与质量数据；GL 扭化把这一切接到 Hitchin 模空间与几何 Langlands。}

